% Source: https://artofproblemsolving.com/community/c3h1881240_polynomial
% Stable identity: aops-topic-1881240-q01
Multiplying the second equation by $pqr$ gives $qr+rp+pq+360=pqr$. Then
\[
(p-1)(q-1)(r-1)=pqr-(pq+qr+rp)+(p+q+r)-1=360+26-1=385=5\cdot7\cdot11 .
\]
Let $a=p-1$, $b=q-1$, $c=r-1$. These are integers with $abc=385$, $a+b+c=23$, and none of them equals $-1$ (since $p,q,r\ne0$).

As $abc>0$, either all three are positive or exactly two are negative.
\begin{itemize}
\item All positive: the factorisations of $385$ into three positive factors are $1\cdot1\cdot385$, $1\cdot5\cdot77$, $1\cdot7\cdot55$, $1\cdot11\cdot35$, $5\cdot7\cdot11$, with sums $387,83,63,47,23$. Only $\{5,7,11\}$ works.
\item Two negative, say $a=-m$, $b=-n$ with $m,n\ge1$: then $c=23+m+n>23$ and $mnc=385$, so $c\in\{35,55,77,385\}$ and $mn=11,7,5,1$ respectively. Then $m+n=12,8,6,2$ and $c-m-n=23,47,71,383$. Only $c=35$, $\{m,n\}=\{1,11\}$ works, but it gives $a$ or $b$ equal to $-1$, which is excluded.
\end{itemize}
So $\{p,q,r\}=\{6,8,12\}$. Check: $6+8+12=26$ and $\dfrac{1}{6}+\dfrac{1}{8}+\dfrac{1}{12}+\dfrac{360}{576}=\dfrac{9}{24}+\dfrac{5}{8}=1$. Hence $pqr=576$ and $\sqrt{pqr}=\boxed{24}$.
